\VerdierPhysicalPageAnchor{0050}%
\VerdierSemanticAnchor{ch1-1-4-7-continuation}
où les $(F_i,p_i)$ sont les $G$-foncteurs tordus déterminés
respectivement par les foncteurs $(F_i,m_i)$ et les sections $s_i$,
sont des $\beta_{j,i}$-morphismes, où :

\VerdierSemanticAnchor{ch1-1-4-7-eq3}
\begin{equation}
  \beta_{j,i}(\sigma)=s_j(\sigma)s_i(\sigma)^{-1},
  \qquad i,j\in\{1,2,3\},\quad i<j.
  \tag{1.4.7.3}
\end{equation}

On a donc la relation :

\VerdierSemanticAnchor{ch1-1-4-7-eq4}
\begin{equation}
  \beta_{3,1}=\beta_{3,2}\beta_{2,1}.
  \tag{1.4.7.4}
\end{equation}

\VerdierSemanticAnchor{ch1-1-4-7-composition}
Les relations (1.4.7.2) et (1.4.7.4) permettent alors de définir le
\emph{composé} du $\beta_{3,2}$-morphisme $v$ avec le
$\beta_{2,1}$-morphisme $u$ et cette composition ainsi définie est
compatible avec la composition des morphismes de $H$-foncteurs.

\VerdierSemanticAnchor{ch1-1-4-8}
\noindent\textbf{1.4.8.}\quad
Soient $\mathcal C'$ une catégorie additive et :
\[
  (F,p)\colon(\mathcal C,G,T,c)\longrightarrow
  (\mathcal C',G,T',c')
\]
un $G$-foncteur tordu par un cocycle normalisé
$\alpha\in Z^2(G,\mathbb Z/2\mathbb Z)$. On désigne par :

\VerdierSemanticAnchor{ch1-1-4-8-eq1}
\begin{equation}
  (F^{\circ},p^{\circ})\colon
  (\mathcal C^{\circ},G^{\circ},T^{\circ},c^{\circ})
  \longrightarrow
  (\mathcal C'^{\circ},G^{\circ},T'^{\circ},c'^{\circ})
  \tag{1.4.8.1}
\end{equation}
le \emph{foncteur gradué tordu opposé} au foncteur $(F,p)$. (On
utilise les notations de (1.3.9)). Il est défini de la manière
suivante :

\VerdierSemanticAnchor{ch1-1-4-8-opposite-data}
Le foncteur sous-jacent est $F^{\circ}$, le foncteur opposé au
foncteur $F$. Pour tout $g\in G^{\circ}$ (groupe opposé à $G$),
l'isomorphisme de $\mathcal C^{\circ}$ :
\[
  p^{\circ}(g)\colon
  F^{\circ}T^{\circ}(g)=FT(g^{-1})
  \overset{\sim}{\longrightarrow}
  T'^{\circ}(g)F^{\circ}=T'(g^{-1})F
\]
est égal à $p^{-1}(g^{-1})$.

\VerdierSemanticAnchor{ch1-1-4-8-opposite-cocycle}
Le foncteur $(F^{\circ},p^{\circ})$ est un $G^{\circ}$-foncteur tordu
par le cocycle $\alpha^{\circ}$ dans
$Z^2(G^{\circ},\mathbb Z/2\mathbb Z)$ opposé au cocycle $\alpha$
\[
  \bigl(\alpha^{\circ}(g,g')
  =\alpha^{-1}(g'^{-1},g^{-1})\bigr).
\]

\VerdierSemanticAnchor{ch1-1-4-9}
\noindent\textbf{1.4.9.}\quad
La seule vertu du groupe $\mathbb Z/2\mathbb Z$ qu'on utilise dans ce
numéro est qu'il opère canoniquement sur le foncteur identique des
catégories additives. On peut donc développer des considérations
analogues en remplaçant le groupe $\mathbb Z/2\mathbb Z$ par un groupe
commutatif $K$ opérant sur les foncteurs identiques des catégories
considérées, de façon compatible avec les opérations de $G$.
